\section{Partecipazione a progetti internazionali}
%
\outerlist{

\entrybig
	{\textbf{NexTArc}}{AimageLab, UNIMORE}
	{\textcolor{darkgray}{\small{Next Generation Open Innovations in Trustworthy Embedded AI Architectures}} [\href{https://cordis.europa.eu/project/id/101194287}{link}]}{09/2025 -- in corso}
\innerlist{
	\entry{\textbf{Attività:} ricerca e sviluppo di tecniche di \textit{Modular Deep Learning} per applicazioni di \textit{smart city} e monitoraggio del traffico nel trasporto pubblico (sistemi di stima e monitoraggio della densità di occupazione dei mezzi di trasporto pubblico, mappatura dei posti disponibili a bordo dei mezzi, etc.). Le attività sono svolte secondo principi di \textit{privacy-preserving} e \textit{trustworthy AI}, garantendo sicurezza e robustezza dei modelli.}
	\entry{\textbf{Finanziamento:} progetto Horizon Europe finanziato dalla Commissione Europea, Grant Agreement n. 101194287.}
}


\entrybig
	{\textbf{STORE}}{AimageLab, UNIMORE}
	{\textcolor{darkgray}{Shared daTabase for Optronic Image Recognition and Evaluation} [\href{https://edf-store.com}{link}]}{12/2023 -- in corso}
\innerlist{
	\entry{\textbf{Attività:} ricerca e sviluppo di sistemi di apprendimento per applicazioni di \textit{object detection}, con particolare riferimento a tecniche di adattamento incrementale dei modelli e \textit{decentralized learning} per l’addestramento e l’aggiornamento di sistemi di riconoscimento visivo in ambito difesa.}
	\entry{\textbf{Finanziamento:} progetto finanziato dalla Commissione Europea nell’ambito dello \textit{European Defence Fund (EDF)}, Grant Agreement n. EDF-2022-101121405.}
}


\entrybig
	{\textbf{InSecTT}}{AimageLab, UNIMORE}
	{\textcolor{darkgray}{Intelligent Secure Trustable Things} [\href{https://www.insectt.eu}{link}]}{2021 -- 2022}
\innerlist{
	\entry{\textbf{Attività:} analisi automatica del comportamento umano nel contesto del trasporto pubblico per l’individuazione di comportamenti anomali, mediante lo sviluppo di modelli di \textit{Deep Learning} e \textit{Computer Vision}.}
	\entry{\textbf{Finanziamento:} progetto europeo Horizon 2020 -- Electronic Component Systems for European Leadership Joint Undertaking (ECSEL JU), Grant Agreement n. 876038.}
}

\entrybig
	{{\textbf{AIDEO}}}{AimageLab, UNIMORE}
	{\textcolor{darkgray}{Artificial Intelligence and Earth Observation data} [\href{https://www.aideo.eu/}{link}]}{2019 -- 2021}
\innerlist{
    \entry{\textbf{Attività:} sviluppo di modelli predittivi basati su \textit{Convolutional Neural Networks} per il monitoraggio della diffusione del virus West Nile, mediante l’analisi di variabili climatiche e ambientali derivate da immagini satellitari multi-banda ad alta risoluzione.}
	\entry{\textbf{Finanziamento:} Agenzia Spaziale Europea (ESA), call \textit{“EO Science for Society – Permanently Open Call for Proposals EOEP-5 Block 4”}.}
}
}